\renewcommand{\thetable}{S\arabic{table}}
\begin{table*}[bt]
\caption{Pollen correction factors from \cite{lang2023quaternary}. Pollen correction factors of 0.25 show overrepresented taxa, while correction factors of 4 show underrepresented taxa. For their usage refer to the equation in the main manuscript.}

\begin{threeparttable}
\begin{tabular}{p{4 cm}  p{5 cm}  p{1.5 cm}}
\headrow
\thead{Group}& \thead{Species}& \thead{Correction factor} \\
Conifers& Abies& 1\\
& Larix& 4\\
& Picea& 1\\
& Pinus Diploxylon (mainly P. sylvestris)& 0.25\\
& Pinus Haploxylon (mainly P. cembra)&1\\
& Juniperus&1\\
Deciduous broadleaved trees& Acer&4\\
& Alnus&0.25\\
& Betula&0.25\\
& Carpinus&1\\
& Castanea&1\\
& Corylus&0.25\\
& Fagus&1\\
& Fraxinus&1\\
& Ostrya&0.25\\
& Quercus robur type&1\\
& Salix&4\\
& Tilia&4\\
& Ulmus&1\\
Evergreen broadleaved trees& Buxus&4\\
 & Quercus ilex type&1\\
& Quercus suber/cerris type&1\\
& Olea&1\\
& Pistacia&1\\
Herbs& Artemisia&1\\
& Chenopodiaceae&1\\
& Plantago lanceolata type&2\\
& Rumex acetosa type&1\\
& Cerealia type&4\\
\end{tabular}
\label{table:2}
\end{threeparttable}
\end{table*}